% \begin{savequote}
% \begin{verbatim}
%  ________________________________________
% / You mentioned your name as if I should \
% | recognize it, but beyond the obvious   |
% | facts that you are a bachelor, a       |
% | solicitor, a freemason, and an         |
% | asthmatic, I know nothing whatever     |
% | about you.                             |
% |                                        |
% | -- Sherlock Holmes, "The Norwood       |
% \ Builder"                               /
%  ----------------------------------------
%        \   ,__,
%         \  (oo)____
%            (__)    )\
%               ||--|| *
% 
% \end{verbatim}
% \end{savequote}

\chapter
[\CO{[TOD]} Sterben und Tod]
{\CC{[TOD]} Sterben und Tod}
\label{chp:TOD}

\ChapterIntro
{%
\noindent Das Sterben und der Tod sind alltägliche Begleiter
des Lebens. Im Gesundheitsbereich ist man zwangsläufig
regelmäßig auf die eine oder andere Weise damit
konfrontiert. Ziel ist es, den legitimen Wünschen der
Patienten und Angehörigen Rechnung zu tragen, und
andererseits selbst mit zum Teil sehr fordernden Situationen
und Fragen zurecht zu kommen.
}
{%
\begin{description}
  \item[Maintainer:] Sebastian Gabriel
  \item[Autoren:] Diverse
  \item[Reviewer:] --
  \item[Version:] {\DocVersion} ({\DocVersionInfo})
  \item[SHA1:] {\DocGitCommit}
\end{description}

{\TxTeilkapitelAASS}
}


\section{Der Tod allgemein}
\index{Tod}
\label{topic:tod}


Der Tod ist eine der grundlegensten und unausweichlichsten
Vitalfunktionen überhaupt.  

\Definition{Der Tod ist das Ende der Lebensfunktionen. Man
unterscheidet zwischen dem klinischen und dem biologischen
Tod}

\paragraph{Klinischer Tod vs. biologischer Tod}
\index{Tod!klinischer}\index{Tod!biologischer}

\begin{center}
\CorrectionItemizeBox\FontTableBase
\begin{tabularx}{\linewidth}{XX}
\toprule
\TabHeading{Klinischer Tod}
 &
\TabHeading{Biologischer Tod}
\\\midrule
keine Atmung

kein Puls

Reanimation muss innerhalb der ersten  Minuten erfolgen

(dann unter Umständen) reversibel
 &
Organtod des Gehirns (irreversibel)

gleichzusetzen mit Tod des Individuums
\\\bottomrule
\end{tabularx}
\end{center}



\section{Begleitung und Betreuung Sterbender}

Dafür gibt es kein Patentrezept und keinen"``goldenen Weg"'.
Wieso? 

Der Tod ist ein integraler Bestandteil des Lebens. So banal
diese Aussage auch scheinen mag, so ambivalent ist das
Verhalten des bzw. der Menschen zu ihm. Nach tausenden
Jahren stellt der Tod und der Sterbeprozess noch immer
(oder vielleicht sogar mehr denn je) eine große
Herausforderung sowohl für den Sterbenden, als auch für
dessen Umgebung dar. 

% \expandedtitle{simple}{\Speech{"`Der Tod ist etwas
% persönliches."'}}

Die Auseinandersetzung mit dem Tod und der eigenen
Endlichkeit verläuft bei jedem Menschen äußerst
unterschiedlich und ist von dessen Weltbild, Umgebung,
Lebensgeschichte und -situation und vielen anderen Faktoren
abhängig, und kann im Verlauf des Lebens ebenso radikalen
Änderungen unterworfen sein. Mögliche Bewältigungstrategien
sind:

\begin{itemize}
  \item Verdrängung
  \item Hoffnungen, welche sich an den Tod anknüpfen
  (Wiedergeburt, Leben im Jenseits)
  \item Vermächtnis an die Nachwelt; Schaffen von etwas, was
 "`größer als man selbst ist"'
  \item Rationalisierung (Vorsorgeversicherungen,
  Testamente, \ldots)
  \item Rituale (Gebete, Begräbniszeremonien,
  Leichenschmaus, \ldots)
\end{itemize}






%%%%%%%%%%%%%%%%%%%%%%%%%%%%%%%%%%%%%%%%%%%%%%%%%%%%%%%%%%%
%%%%%%%%%%%%%%%%%%%%%%%%%%%%%%%%%%%%%%%%%%%%%%%%%%%%%%%%%%%
%%%%%%%%%%%%%%%%%%%%%%%%%%%%%%%%%%%%%%%%%%%%%%%%%%%%%%%%%%%
%%%%%%%%%%%%%%%%%%%%%%%%%%%%%%%%%%%%%%%%%%%%%%%%%%%%%%%%%%%
\subsection{Der todkranke und sterbende Patient}

Rettungsdienst: Was erwartet der Patient und dessen
Angehörigen?

% \expandedtitle{roundbox}{\Speech{"`Das große Problem hat das
% Fachpersonal."'}}





%%%%%%%%%%%%%%%%%%%%%%%%%%%%%%%%%%%%%%%%%%%%%%%%%%%%%%%%%%%
%%%%%%%%%%%%%%%%%%%%%%%%%%%%%%%%%%%%%%%%%%%%%%%%%%%%%%%%%%%
%%%%%%%%%%%%%%%%%%%%%%%%%%%%%%%%%%%%%%%%%%%%%%%%%%%%%%%%%%%
%%%%%%%%%%%%%%%%%%%%%%%%%%%%%%%%%%%%%%%%%%%%%%%%%%%%%%%%%%%
\subsection{Der unerwartete Tod}
% \begin{itemize}
%   \item Das Problem des Patienten bleibt grundsätzlich das
%   Problem des Patienten
%   \item Der Helfer muss im Rahmen der Möglichkeiten im Sinne
%   des Patienten handeln und seine angemessenen
%   Bedürfnisse befriedigen.
%   \item Wenn der Helfer mit unzufireden mit der Betreuung
%   ist, muss er versuchen, den Rahmen des Möglichen
%   zu erweitern.
% \end{itemize}
% 
% \begin{itemize}
%   \item Man darf nicht seine eigenen Vorstellungen
%   bezüglich der Hilfe mit denen des Patienten
%   gleichstellen. Wenn man dem Patienten nutzen möchte, muss
%   man wissen, was \E{er} tatsächlich will.
% \end{itemize}



\section{Todesfeststellung}


\AuthNotes{Todeszeichen: Nicht wie von Roman gefordert in
Tabelle mit mit sicheren und unsicheren Todeszeichen
umgebaut, halte ich für unübersichtlich und befürchte
Verwechslungsgefahr mit Knochenbruchzeichen.[GAB]}





\Layout{2}{Zeichen des Todes und Leichenzeichen}
{%
Grundsätzlich gibt es 3 Zeichen die als \EE{sichere
Todeszeichen} gelten~\cite{ForMed2}:

\begin{enumerate}
    \item \EE{Bildung von blauvioletten Totenflecken} durch
    Absackung des Blutes an die untenliegenden Stellen.
    \item \EE{Nicht mit dem Leben vereinbare Verletzungen}
    wie zum Beispiel eine Abtrennung des Kopfes vom Körper.
    \item \EE{Fäulniserscheinungen}
\end{enumerate}

\noindent Die ersten beiden Punkte sind \E{frühe Leichenerscheinungen}, da sie
relativ rasch auftreten. Andere frühe Leichenerscheinungen sind die \E{Abkühlung der
Leiche} und die \E{Totenstarre}. Längere Zeit nach dem Tod treten die \E{späten
Leichenerscheinungen} auf, dazu gehören \E{Fäulnis und Verwesung},
\E{Mumifizierung}, \E{Tierfraß} und \E{Skelettierung}.

Die Feststellung der Todeszeichen und Leichenerscheinungen
ist jedoch oft nicht so einfach, die Anzeichen können oft
auch mit anderen Sachen verwechselt werden, z.\,B.
Totenflecke mit Blutergüssen oder die Totenstarre mit einem
tonischen Krampf.
}{}{}{}{}


\Fazit{Die Todesfeststellung erfolgt grundsätzlich nur durch einen
Arzt!}\index{Todesfeststellung}

\Cave{Nobody is dead until warm and dead!

Niemand wird für tot erklärt, bevor er nicht erwärmt
wurde.}

%\vspace{1em}
%%%%%%%%%%%%%%%%%%%%%%%%%%%%%%%%%%%%%%%%%%%%%%%%%%%%%%%%%%%
%%%%%%%%%%%%%%%%%%%%%%%%%%%%%%%%%%%%%%%%%%%%%%%%%%%%%%%%%%%
%%%%%%%%%%%%%%%%%%%%%%%%%%%%%%%%%%%%%%%%%%%%%%%%%%%%%%%%%%%
\section{Unterlassung der Reanimation}
\label{topic:reanimation-unterlassung}




\MassnahmenWide{Spezielle Maßnahmen}{Unterlassung der Reanimation}
{\label{m:reanimation-unterlassung}}
{%
\fcolorbox{ubuntured}{White}{
\begin{minipage}[t]{\linewidth - {4\fboxrule} - {4\fboxsep}}
\Ca{Unter den folgenden Voraussetzungen darf (oder muss) eine
Reanimation unterlassen
werden:\index{Reanimation!Unterlassung
der}\index{Wiederbelebung|see{Reanimation}}}
\begin{multicols}{2}              
\begin{itemize}
  \item \EE{Verwesung}
  \item \EE{Fäulnis}
  \item \EE{Mumifizierung}
  \item \EE{Skelettierung}
  \item \EE{Verletzungen, die nicht mit dem Leben
  vereinbar sind}
  \item \EE{Tierfraß}\footnote{Unter
  \E{"`Tierfraß"'} versteht man den Verzehr eines Kadavers
  durch Kadaverfresser. Ein frischer Hundebiss oder
  Ähnliches fällt \underbar{nicht} darunter. Maden können
  sich im Gewebe von Lebenden einnisten und sind daher
  auch kein sicheres Todeszeichen.}
  \item (Verbindliche) \EE{Patientenverfügung}\footnote{\E{Patientenverfügung}: Bei Notfallversorgungen darf \EE{nicht nach einer Verfügung gesucht
  werden, wenn dadurch das Leben oder die Gesundheit des Patienten ernstlich
  gefährdet werden würde}.
  \EE{Aber:} Wenn \E{zweifelsfrei} eine verbindliche
  Patientenverfügung vorliegt (z.\,B. weil eine betreuende
  Pflegekraft diese in der Zwischenzeit gefunden und die
  Verbindlichkeit bestätigt hat), muss dieser Folge
  geleistet werden.
Liegt eine \E{beachtliche Patientenverfügung} vor, so ist im
\E{Einzelfall} zu entscheiden: Das Wohl des Patienten bleibt
oberstes Gebot und die Verfügung soll in die
Entscheidung einfließen.}\index{Patientenverfügung!Unterlassung der
  Reanimation} (\FullPrettyRef{topic:patientenverfuegung}; 
  Bis zur zweifelsfreien Klärung der Situation muss
  jedenfalls eine Reanimation begonnen bzw. fortgeführt
  werden!)
  \item (Abbruch der Reanimationsmaßnahmen aufgrund
  von \EE{Erschöpfung})\footnote{Erschöpfung darf im
  professionellen Umfeld keine Rolle spielen, für eine
  entsprechende Ablösung ist bei Bedarf zu sorgen.}
  \item Unvereinbarkeit mit
  \EE{Selbstschutz}\footnote{Es müssen jedoch alle Maßnahmen
  getroffen werden, die ohne Gefährdung des Selbstschutzes
  möglich sind. So wird zum Beispiel für Ersthelfer
  empfohlen, wenn
  eine Mund-zu-Mund-Beatmung nicht zumutbar ist, zumindest 
  eine Herzdruckmassage durchzuführen.}
\end{itemize}

\end{multicols}
\end{minipage}
}
}








\AuthNotes{GABG: gesicherte bösartige Erkrankung?

GABS: Was gilt als gesichert? Was gilt als bösartig? Eine
bösartige Grunderkrankung gilt nicht per se als Argument um
nicht zu reanimieren $\rightarrow$ Patientenverfügung,
mutmaßlicher Patientenwille. Prostatakarzinom und Infarkt?}







%O2:
%\nocite{Morris2001,Milross1989,Boumphrey2003}
%\nocite{BoehmerTaschRett6,Fauci2008,LPN3,LippertAnatomie7,StriebelAn7,DRAn3}


%\putbib

\ChapterBibliography



